\section{Refinement and its guarantees}\label{sec:legal}

We now prove the refinement guarantees used in Section~\ref{sec:potential}:
moving an item preserves the answers already given, and the resulting
change in volume controls the cost of choosing a child.

\subsection{Preserving the represented orders}

We use the interpretation of the state given in Section~\ref{sec:overview}.
Assignments with equal free coordinates, or with a free coordinate equal
to an anchor, lie on a finite union of sets of volume zero and are excluded
when reading an order. We maintain the following invariant:
\begin{quote}
Every represented total order respects every comparison already answered.
\end{quote}
The state may discard some orders that still respect the answers. Legality
requires only that at least one remains. Positive volume in every node
will ensure this.

The invariant must continue to hold when refinement removes an arc from
an in-node DAG.

\begin{lemma}[Effective intervals and sections]\label{lem:interior}
Replacing interval bounds by their effective endpoints leaves the
feasible region unchanged.
The region has positive volume if and only if every effective width is
positive. When these widths are positive, fixing $X_x=c$ at any
$c\in(\ell_x,r_x)$ leaves a
positive-volume section in the remaining coordinates.
\end{lemma}
\begin{proof}
If $y\ne x$ precedes $x$ on a directed path, every feasible assignment
satisfies $X_x>X_y>a_y$; the stored bound $X_x>a_x$ covers $y=x$.
Taking the maximum over these bounds gives $X_x>\ell_x$. The upper-bound argument is the reverse. Thus all bounds
given by \eqref{eq:effective} already hold.

Positive volume requires positive widths. Conversely, suppose all widths
are positive. Along an arc $x\to y$, the definitions give
$\ell_x\le\ell_y$ and $r_x\le r_y$. Choose numbers $0<\theta_x<1$
increasing in a topological order, and set
\[
X_x=(1-\theta_x)\ell_x+\theta_xr_x.
\]
Increasing $\theta_x$ to $\theta_y$ strictly increases this value because
$r_x>\ell_x$. Replacing its endpoints by the larger endpoints
$\ell_y,r_y$ cannot decrease it. Hence $X_x<X_y$. All coordinates are
inside their intervals and all arc inequalities hold strictly, so a
sufficiently small open neighborhood of this assignment is feasible.

Now fix $X_x=c$ and delete its coordinate. The inequalities imposed by
its incident arcs become upper bounds $c$ at its in-neighbors and lower
bounds $c$ at its out-neighbors. Upper bounds propagate only backward
along arcs, and lower bounds only forward. Any lower bound that previously
reached a successor through $x$ is superseded by $c>\ell_x$; any upper
bound that reached a predecessor through $x$ is superseded by $c<r_x$.
Thus the new effective intervals are
\[
\begin{cases}
(\ell_y,\min(r_y,c)),&y\text{ precedes }x,\\
(\max(\ell_y,c),r_y),&y\text{ succeeds }x,\\
(\ell_y,r_y),&y\text{ is incomparable with }x.
\end{cases}
\]
The first two cases are disjoint because the in-node DAG is acyclic. A predecessor
retains positive width: $r_y>\ell_y$ and
$c>\ell_x\ge\ell_y$. A successor does too: $r_y>\ell_y$ and
$c<r_x\le r_y$. All other widths stay positive. The criterion just proved,
applied to the remaining in-node DAG, shows that the section has positive volume.
\end{proof}

\begin{lemma}[Refinement and child admission]\label{lem:admission}
Suppose the node regions have positive volume and the order invariant holds.
Refining $x$ at an interior anchor $c$ preserves both properties, including
when the child at $c$ already exists. The transcript implies no
relation between $x$ and any label already in that child subtree.
\end{lemma}
\begin{proof}
Consider first an existing child at $c$, and denote its label set by $S$.
Choose a feasible assignment in the parent section $X_x=c$, with the other
free coordinates pairwise distinct and different from every child anchor.
Such an assignment exists by Lemma~\ref{lem:interior}. Also choose feasible
assignments in the other nodes.

There is a positive distance from $c$ to every other child anchor and
every free parent coordinate in this assignment. All incident inequalities at $x$ are strict. Thus, with
everything else fixed, $x$ can be moved slightly to either side of $c$
while keeping the old state feasible. The resulting two represented orders
place $x$ immediately before or immediately after the entire block $S$.
Both orders respect the transcript. For each $y\in S$, one has $x<y$ and
the other has $y<x$; neither relation can therefore be implied by the
transcript.

Now take any order represented after the refinement. Deleting $x$ from the
child leaves an assignment of its old coordinates and hence an old child
order. Apply the preceding perturbation using this assignment's parent
coordinates and old child coordinates. It gives two feasible old orders
that agree with the new one on every label except $x$. The new order
inserts $x$ somewhere into their child block. No
answered comparison joins $x$ to a label in $S$, so this insertion violates
none of them. Relative to every label outside $S\cup\{x\}$, the whole
block occupies the old anchor position. The item $x$ is consequently on
the same side of each outside label as in either of the two feasible old
orders above. Comparisons between other labels are unchanged. The new
order therefore respects the entire transcript.

If there was no child at $c$, the same argument uses an empty block: the
new singleton child places $x$ where its fixed parent coordinate stood.
Finally, the remaining parent region is a positive-volume section, and
the child receives an independent coordinate in $(0,1)$. All node regions
remain positive.
\end{proof}

\subsection{The volume gained by refinement}\label{sec:refinement-volume}

Refining an item $x$ stored at node $v$ fixes $X_x=c$ in $v$ and gives
$x$ a fresh unit coordinate in the child. To measure the volume change, let
$S_x(c)$ be the volume of the section of $K_v$ obtained by fixing
$X_x=c$, measured in the remaining coordinates. Under the uniform
distribution on $K_v$, the coordinate $X_x$ has density
\[
f_x(c)=\frac{S_x(c)}{V_v},\qquad
\Prb(a<X_x<b)=\int_a^b f_x(t)\,dt.
\]

\begin{lemma}[Refinement identity]\label{lem:section-volume}
Suppose every node region has positive volume. Refining $x$ at an interior
anchor $c$ satisfies
\begin{equation}\label{eq:refine}
\frac{Z_{\rm after}}{Z_{\rm before}}=f_x(c).
\end{equation}
\end{lemma}
\begin{proof}
Let $u$ be the child at $c$, taking $V_u=1$ if it is newly created. The
parent loses the coordinate $X_x$ and retains precisely the section of
volume $S_x(c)$. The child gains a new coordinate $U_x\in(0,1)$, with no
constraints involving its old coordinates. Its region becomes
$K_u\times(0,1)$ and still has volume $V_u$. Therefore the affected
factors in the product change as follows:
\[
V_vV_u\quad\longrightarrow\quad S_x(c)V_u.
\]
Their ratio is $S_x(c)/V_v=f_x(c)$. All other node factors stay unchanged.
The dimension lost in the parent is restored by the new child coordinate.
\end{proof}

Thus the refinement's volume gain is $g=\log f_x(c)$.

For example, the parent region $0<X_A<X_B<1$ has area $1/2$. Fixing
$X_A=1/4$ leaves the interval $X_B\in(1/4,1)$ of length $3/4$. The new
child coordinate $U_A\in(0,1)$ turns this into a rectangle of area $3/4$,
so refinement multiplies $Z$ by $3/2$; see Figure~\ref{fig:section-volume}.
Both regions represent the order $A<B$. Replacing the fixed parent
coordinate by a fresh child coordinate gives the larger area.

\begin{figure}[tb]
\centering
\begin{tikzpicture}[x=3.25cm,y=3.25cm,>=Latex,font=\small]
\begin{scope}
\node at (.5,1.38) {Before refinement};
\fill[pale] (0,0)--(0,1)--(1,1)--cycle;
\draw[ink!55] (0,0)--(1,1)--(0,1);
\draw[->] (0,0)--(1.12,0) node[right] {$X_A$};
\draw[->] (0,0)--(0,1.12) node[above] {$X_B$};
\draw[dashed,ink!60] (.25,0)--(.25,.25);
\draw[accent,very thick] (.25,.25)--(.25,1);
\node[below] at (.25,0) {$\frac14$};
\node[below left] at (0,0) {$0$};
\node[below] at (1,0) {$1$};
\node[left] at (0,1) {$1$};
\node at (.55,.83) {area $\frac12$};
\end{scope}
\draw[->,thick] (1.2,.53)--(1.63,.53);
\node[above] at (1.415,.56) {refine $A$};
\begin{scope}[shift={(1.9,0)}]
\node at (.5,1.38) {After refinement};
\fill[pale] (0,.25) rectangle (1,1);
\draw[ink!55] (0,.25)--(1,.25)--(1,1)--(0,1);
\draw[->] (0,0)--(1.12,0) node[right] {$U_A$};
\draw[->] (0,0)--(0,1.12) node[above] {$X_B$};
\node[below left] at (0,0) {$0$};
\node[left] at (0,.25) {$\frac14$};
\node[below] at (1,0) {$1$};
\node[left] at (0,1) {$1$};
\node at (.5,.65) {area $\frac34$};
\end{scope}
\end{tikzpicture}
\caption{The highlighted section of the triangle has length $3/4$.
Replacing the fixed parent coordinate by a fresh unit coordinate in the
child gives the rectangle on the right. Its area is $3/2$ times the
original area.}
\label{fig:section-volume}
\end{figure}

\subsection{Choosing an anchor}\label{sec:median-search}\label{sec:anchor-contract}

The choice of anchor balances two quantities: a larger density gives
more volume gain, while a smaller anchor depth costs less in
\eqref{eq:anchor-cost}. We will bound the density using the number of
coordinates comparable with $x$. Put
\begin{equation}\label{eq:comparable-bound}
B=(k+1)^{2H+1}.
\end{equation}
In an in-node DAG of degree at most $k$ and directed paths of at most
$2H$ vertices, every coordinate comparable with $x$ lies within undirected
distance $2H$ of $x$. Thus $B$ bounds the number of other coordinates
comparable with $x$.

Set the calibration width
\begin{equation}\label{eq:width-threshold}
w_*=2^{-2kB}.
\end{equation}
The finite calculations use this scale to guarantee success when a
queried coordinate has a sufficiently large effective interval.

\begin{lemma}[When certification succeeds]\label{lem:certificate-success}
Suppose a node has positive feasible volume, its in-node DAG has degree
at most $k$, and its directed paths have at most $2H$ vertices.
If any free coordinate involved in a comparison has effective width at
least $w_*$, comparison certification succeeds. If $W_x\ge w_*$, then
\textnormal{\texttt{CentralInterval}}$(v,x)$ succeeds.
\end{lemma}

Section~\ref{sec:oracle} proves this lemma. In particular, failure of a
central-interval certificate implies $W_x<w_*$. This is the information
used to justify the anchor routine's fallback.

\begin{lemma}[Anchor selection]\label{lem:anchor-density}
Suppose a node's feasible region has positive volume, its interval
endpoints are dyadic, and its in-node DAG has degree at most $k$ and
directed paths of at most $2H$ vertices. With calibration width \eqref{eq:width-threshold},
\textnormal{\texttt{ChooseAnchor}} deterministically returns an anchor
$c\in(\ell_x,r_x)$ at depth $m$ in the implicit anchor tree satisfying
\begin{equation}\label{eq:tariff}
\log f_x(c)\ge(1-1/k)m-4.
\end{equation}
If $x$ is refined because comparison certification failed, then
\begin{equation}\label{eq:failure-gain}
\log f_x(c)>k.
\end{equation}
\end{lemma}

Algorithm~\ref{alg:choose-anchor} first tries \texttt{CentralInterval}.
On success, it uses the returned interval; on failure, it uses the middle
half of the effective interval. It then selects the shallowest anchor in
the chosen interval. Both choices give the density bound required by
Lemma~\ref{lem:anchor-density}.

\begin{algorithm}[H]
\caption{\texttt{ChooseAnchor}$(x)$: choose a position for refinement.}
\label{alg:choose-anchor}
\DontPrintSemicolon
\KwIn{A free item $x$ stored at node $v$ and its effective interval $(\ell_x,r_x)$.}
\KwOut{A dyadic anchor $c\in(\ell_x,r_x)$.}
$J\gets$ \texttt{CentralInterval}$(v,x)$\;
\If{$J$ is \textnormal{\texttt{FAIL}}}{
  $W_x\gets r_x-\ell_x$\;
  $J\gets[\ell_x+W_x/4,\ r_x-W_x/4]$\;
}
$[a,b]\gets J$\;
$u\gets$ the root of the implicit anchor tree\;
\While{true}{
  \uIf{$c(u)<a$}{$u\gets$ the right child of $u$\;}
  \uElseIf{$c(u)>b$}{$u\gets$ the left child of $u$\;}
  \Else{\Return $c(u)$\;}
}
\end{algorithm}

An anchor at depth $m$ has $m+1$ fractional binary digits in reduced
form. Thus the walk finds the shortest binary fraction in $[a,b]$.

\begin{lemma}[Shortest anchor in an interval]\label{lem:shortest-anchor}
For a nondegenerate closed interval $[a,b]\subset(0,1)$, the walk in
Algorithm~\ref{alg:choose-anchor} returns its unique minimum-depth anchor.
If its node is $u$ and its depth is $m$, then
$[a,b]\subseteq I(u)$ and $|I(u)|=2^{-m}$.
\end{lemma}
\begin{proof}
Initially $[a,b]\subseteq I(u)=[0,1]$. Whenever $c(u)<a$, the whole
interval lies in the right child's interval; whenever $c(u)>b$, it lies
in the left child's interval. The rejected anchor and the other subtree
contain no candidate in $[a,b]$. The first accepted anchor therefore has
minimum depth, and every other anchor of that depth lies outside $I(u)$.
The intervals halve at every step and continue to contain $[a,b]$, so
the walk must accept before their length becomes smaller than $b-a$.
\end{proof}

\subsubsection{Density bounds for central intervals}

The first bounds use the effective width and the number of comparable
coordinates.

\begin{lemma}[Coordinate density bounds]\label{lem:fiber}
Suppose a node's feasible region has positive volume, and $x$ has at most
$B$ other comparable coordinates and effective interval $(\ell,r)$ of
width $W$. Its density satisfies
\begin{align}
f_x(t)&\le\frac{B+1}{W},\label{eq:fiber-upper}\\
f_x(t)&\ge\frac{4^{-B}}{W}
\qquad(\ell+W/4\le t\le r-W/4).\label{eq:fiber-lower}
\end{align}
\end{lemma}

The lower bound controls the fallback choice in the middle half of the
effective interval. The upper bound will help bound the probability of
an unresolved comparison. Appendix~\ref{app:geometry} proves both by
comparing convex sections.

For a certified central interval we use another consequence of convexity. A density is
\emph{log-concave} if its positive support is convex and its logarithm
is concave there. The uniform density on a convex feasible region is
log-concave, and Pr\'ekopa's marginal theorem gives the same property to
each coordinate density \cite[Theorem~8]{prekopa}. The following bound says
that a position with substantial probability on each side has density
comparable to the largest value.

\begin{lemma}[Density at central positions]\label{lem:central-quantiles}
Let $X$ have a log-concave density $f$. If
$\Prb(X\le c)\ge p$ and $\Prb(X\ge c)\ge p$ for some $0<p\le1/2$, then
\[
f(c)\ge p\sup_t f(t).
\]
\end{lemma}

Appendix~\ref{app:geometry} proves the lemma directly from log-concavity.

\subsubsection{Anchors from a certified interval}

\begin{lemma}[Certified anchor choices]\label{lem:median-search}
If \textnormal{\texttt{CentralInterval}} succeeds, the anchor $c$ selected
by Algorithm~\ref{alg:choose-anchor} has depth $m$ and satisfies
\[
\log f_x(c)\ge m-4,\qquad
\log f_x(c)\ge\log(1/W_x)-2.
\]
\end{lemma}
\begin{proof}
Let $J$ be the interval returned by \texttt{CentralInterval}, and put
$M=\sup_t f_x(t)$. By Lemma~\ref{lem:central-interval}, $J$ is contained
in $(\ell_x,r_x)$ and has probability at least $1/4$, so $M|J|\ge1/4$.
Every position $c\in J$ also has probability at least $1/4$ on each side.
Lemma~\ref{lem:central-quantiles} therefore gives $f_x(c)\ge M/4$
throughout $J$.

Lemma~\ref{lem:shortest-anchor} selects $c\in J$ at a node $u$ whose
interval contains $J$. Its interval has length $2^{-m}$, so
\[
2^{-m}f_x(c)\ge |J|\frac M4\ge\frac1{16}.
\]
This proves the first bound. Since $f_x$ integrates to one over an
interval of width $W_x$, its supremum satisfies $M\ge1/W_x$.
The inequality $f_x(c)\ge M/4$ proves the second bound.
\end{proof}

\subsubsection{The anchor guarantee}\label{sec:anchor-guarantee}

The certified choice already gives the required bound. For the fallback,
failure of the central-interval calculation implies $W_x<w_*$, so
$\log(1/W_x)>2kB$. This pays for the loss of $2B$ in the lower density
bound \eqref{eq:fiber-lower}.

\begin{proof}[Proof of Lemma~\ref{lem:anchor-density}]
The in-node DAG bounds imply that $x$ has at most $B$ other comparable
coordinates, by \eqref{eq:comparable-bound}. If central certification
succeeds, Lemma~\ref{lem:median-search} gives the stronger lower bound
$m-4$. Otherwise, Lemma~\ref{lem:certificate-success} gives
$W_x<w_*$. Write $L=\log(1/W_x)>2kB$. The fallback anchor is in the
central half of its effective interval, so \eqref{eq:fiber-lower}
gives
\[
\log f_x(c)\ge L-2B.
\]
By Lemma~\ref{lem:shortest-anchor}, its node interval contains that central
half. Hence $W_x/2\le2^{-m}$ and $m\le L+1$, giving
\[
\log f_x(c)\ge L-2B>(1-1/k)L
\ge(1-1/k)(m-1)>(1-1/k)m-4.
\]
This proves \eqref{eq:tariff} in both branches.

Now suppose comparison certification failed and the algorithm selected
$x$ for refinement. Lemma~\ref{lem:certificate-success} implies
$W_x<w_*$, regardless of the subsequent central-interval result.
If central certification succeeds, Lemma~\ref{lem:median-search} gives
\[
\log f_x(c)\ge\log(1/W_x)-2>2kB-2>k.
\]
If it fails, the fallback bound above gives
$\log f_x(c)\ge\log(1/W_x)-2B>2B(k-1)>k$.
This proves \eqref{eq:failure-gain}.
Both cases use a finite deterministic calculation to obtain the central
interval, followed by the finite walk of Lemma~\ref{lem:shortest-anchor}.
\end{proof}

\subsection{Bounds during a comparison call}\label{sec:repairs}

The preparation loop in Algorithm~\ref{alg:compare} simplifies the current
node before an answer is selected. We check that the in-node DAG bounds
hold throughout the call, and that the loop makes only finitely many
moves.

\begin{lemma}[In-node DAG bounds and termination]\label{lem:restoration}
Suppose every node region has positive volume and every in-node DAG has
degree at most $k$ and directed paths of fewer than $2H$ vertices at the
start of a call. These properties hold throughout the call. Each label
is refined at most once, so the call makes at most $n$ refinements.
\end{lemma}
\begin{proof}
A refinement uses an interior anchor by Lemma~\ref{lem:anchor-density}.
It leaves a positive-volume parent section by Lemma~\ref{lem:interior}
and gives the child a fresh unit coordinate. It deletes a vertex and its
incident arcs from the parent and adds an isolated vertex to the child,
so it cannot increase any degree or extend any directed path.

A returned comparison certificate gives its chosen answer positive
probability, so recording it preserves positive feasible volume.
Its arc cannot create a directed cycle, which would make the region
empty. Immediately before a new answer is selected,
the current in-node DAG has no path on $H$ vertices and the queried free
vertices have degree below $k$. Adding their comparison arc leaves degree
at most $k$. Any new directed path consists of two old paths joined by
this arc, so it has fewer than $2H$ vertices. An answer involving an anchor
adds only an interval bound. These observations prove positivity and
the in-node DAG bounds by induction through the call.

Once a label is refined, it has an isolated coordinate with interval
$(0,1)$ in its child. Until the final answer, every subsequent state change
is another refinement. These moves affect interval bounds only at
neighbors of the vertex being removed, so they cannot narrow this isolated
coordinate or give it an incident arc. It cannot lie on a selected path
or have degree $k$. Lemma~\ref{lem:estimation} guarantees that comparison
certification succeeds whenever such a unit-interval coordinate is
queried, so it cannot be selected after a failed certificate either.
Thus no label is selected again, even if the query is rerouted to its new
node. Every restart refines at least one label, giving the bound of $n$
refinements.
\end{proof}

\begin{proof}[Proof of Proposition~\ref{prop:legality}]
Initially every node region has positive volume, every order is represented,
and the in-node DAG bounds hold. We induct over the changes made by a
comparison call. The bounds follow from Lemma~\ref{lem:restoration}.
At a refinement, Lemma~\ref{lem:anchor-density} supplies an interior
anchor, and Lemma~\ref{lem:admission} preserves positive volume and the
order invariant. This applies to every item of a selected path, including
those whose incident arcs have already been removed by earlier moves.

If the queried labels lie in different child subtrees, their anchors
determine the answer in every represented order. Otherwise, a comparison
certificate supplies a chosen answer retaining at least $1/2-\eta>0$
of the current node's volume, by the acceptance test in
Algorithm~\ref{alg:compare} and Lemma~\ref{lem:estimation}.
Recording its arc or interval bound makes that answer true in every
represented order and preserves every earlier constraint. This completes
the induction.

The comparison calculation and anchor search terminate by
Lemmas~\ref{lem:estimation} and~\ref{lem:anchor-density}. Together with
Lemma~\ref{lem:restoration}, this proves termination of the full call.
These analytic lemmas apply to arbitrary positive-volume regions with
the stated in-node DAG bounds; their proofs are independent of the
adversary's history.
\end{proof}
